\documentclass[11pt]{article}
\usepackage[margin=1.1in]{geometry}
\usepackage{amsmath,amssymb,amsthm}
\usepackage{booktabs}
\usepackage{array}
\usepackage{tabularx}
\usepackage[protrusion=true,expansion=false]{microtype}
\usepackage[table]{xcolor}
\usepackage{enumitem}
\usepackage{titlesec}
\usepackage{float}
\usepackage{hyperref}

\hypersetup{colorlinks=true, linkcolor=blue!60!black, urlcolor=blue!60!black, citecolor=blue!60!black}
\usepackage[skip=0.6\baselineskip, indent=0pt]{parskip}

\titleformat{\section}{\large\bfseries}{\thesection.}{0.7em}{}
\titleformat{\subsection}{\normalsize\bfseries}{\thesubsection.}{0.6em}{}

\newcommand{\vmax}{V_{\max}}
\newcommand{\Tc}{T_{\mathrm{close}}}
\newcommand{\Astar}{A^{*}}
\newcommand{\tx}[1]{\texttt{#1}}
\newcolumntype{L}{>{\raggedright\arraybackslash}X}

\renewcommand{\arraystretch}{1.3}
\makeatletter
\setlength{\@fptop}{0pt}
\setlength{\@fpsep}{1.5\baselineskip}
\setlength{\@fpbot}{0pt plus 1fil}
\makeatother
\rowcolors{1}{}{black!5}

\title{\textbf{Lichen's reference L2}\\as built and measured\\[2pt]\large draft}
\author{waxwing}
\date{Version 0.1, November 2026\\[2pt]\small a companion to \emph{Lichen}, version 0.5~\cite{lichen}}

\begin{document}
\maketitle

\begin{abstract}
Lichen~\cite{lichen} enforces an L2's withdrawals through bilateral
contracts between a hub and each user, settled on Bitcoin by BABE: a
user claims, and only the hub proves, if it can, that her claim is
false. This companion describes the L2 that lichen builds to show it,
in which the hub is also the sequencer, so that everything a separate
L2 would vouch for must be enforced against the operator. Its ledger is
one of notes, transparent and shielded, in which every block is proved
valid by a chain of RISC Zero proofs and the history is anchored on
Bitcoin by a chain of single-use seals. A spend counts against its
sender only once she has acknowledged it against an anchor, with its
data in hand, so withholding data stalls payments but cannot take
funds. An exit claims a definite set of notes, judged at the last anchor
a margin before the claim; the hub's statement refutes it on eight
grounds, all shown from data the operator holds. We give the design and
measurements from a deployment on Mutinynet (November 2026) at real
parameters: an honest exit is two transactions and 2.7~kvB with no
prover on the user's side; a stale claim was proved false on the hub's
16-vCPU server in about 9 minutes and wrapped in 2, and in 3.5~s on a
rented GPU; a payment is firm in 10 to 20 minutes on Mutinynet, about an
hour on mainnet. The security argument is a sketch, and the hub's
proving capacity against a wave of false claims is the main open
question.
\end{abstract}

\tableofcontents

\section{Scope}
\label{sec:scope}

The main paper~\cite{lichen} describes a contract between a user,
Alice, and a hub: a 2-of-2 output holding her deposit $D$, the hub's
liquidity up to a cap $\vmax$, and two reserves $r_A, r_H$. The hub's
tokens on an L2 are collateralised by it. A withdrawal is a claim $C$
of $b$, folded cooperatively, or, if the hub refuses or is offline,
posted on Bitcoin in \tx{request}, where the hub must prove it false or
let a timeout pay her. The hub's statement is the L2's:
\[
  \mathcal{R}_H\big(\mathcal{D}, c, \ldots;\; w = (\mathsf{req}, \rho,
  \ldots)\big) \;=\; \neg\,\mathsf{Valid}_L\big(C(\mathsf{req}),\,
  b(\mathsf{req})\big),
\]
with $\rho$ a refutation. How BABE checks the hub's proof on Bitcoin,
with the roles swapped so that Alice garbles and the hub proves, and the
setup by verifiable secret sharing, are in the main paper's
appendix~\cite[App.~A]{lichen}. This companion is the other side of the
interface: lichen's own L2~\cite{lichencode}, what $\mathsf{Valid}_L$
and its refutations are for it, and how it all measured.

It replaces \emph{Lichen as built}~\cite{lichenbuilt}, whose two parts
described an earlier design: Alice proved her own withdrawal, the hub
garbled, a withdrawal began with a return of tokens to the hub, and an
exit could be answered only with a counter that Alice might then have
to prove false. All of that is gone (Section~\ref{sec:gone}).

Figures marked (lichen) are measured: on Mutinynet, a signet with
30-second blocks, on one deployment opened on 2026-11-11 with seven
contracts at real parameters, or on regtest where stated. The rest are
estimates, marked (est.). The machines are a 14-core Apple M4 Pro
laptop (RISC Zero has no Metal prover, so its CPU); \emph{lichen-hub},
the deployment's own server, 16 x86 vCPUs at 2.0~GHz without AVX; a
4-vCPU x86 server; and a rented NVIDIA RTX 4090.
A walk-through of every Mutinynet run, transaction by transaction, is
in lichen's explainers~\cite{explainers}.

\section{When the hub is the sequencer}
\label{sec:merged}

When the L2 is not the hub's (an L2 secured by its own consensus, or a
rollup elsewhere with its own proofs), the L2's account of its ledger is
an input, and the hub only a counterparty. In lichen's L2 the hub and
the sequencer are one operator, or two that trust each other. That is a
design choice, and its consequence is that the L2's account is the
operator's own word. A hub profits from defaults, since whatever isn't
withdrawn by $\Tc$ is its, so it has every reason to bend that account,
and the hub's statement must not take it on trust.

With hub and sequencer one party, the threat is the operator against its
users. Users against users is the ledger's ordinary job, and the
operator against itself is out of scope. Table~\ref{tab:threats} lists
what the operator could do, and what answers it.
\begin{table}[!htbp]
\centering
\small
\begin{tabularx}{\linewidth}{>{\raggedright\arraybackslash}p{62mm}L}
\toprule
the operator & answered by \\
\midrule
mints notes nobody paid for, and uses them to refute a claim & validity: the chain proof (Section~\ref{sec:validity}) \\
signs two valid histories that split & one anchored history: the seal chain (Section~\ref{sec:finality}) \\
anchors a history on a branch of Bitcoin it mined & the contract's work $W$ (Section~\ref{sec:finality}) \\
anchors a spend of Alice's and withholds its data & sender acknowledgment: a spend counts only once she acknowledges it, with her data (Section~\ref{sec:ack}) \\
holds a payment of hers unsequenced, then includes it once she claims its inputs & the margin: a claim is judged at an anchor fixed before it (Section~\ref{sec:claim}) \\
refuses the fold, censors her, stops sealing, or vanishes & the exit: her claim, which the hub must prove false (Sections~\ref{sec:claim}--\ref{sec:statement}) \\
claims her notes through a contract of its own & only her keys sign a claim of her notes \\
\bottomrule
\end{tabularx}
\caption{What a merged operator could do, and what answers it.}
\label{tab:threats}
\end{table}

\section{The ledger}
\label{sec:ledger}

The ledger is one of notes. A \emph{transparent} note is an owner key, a
value and a memo, with a random $\rho$; a \emph{shielded} note is a
commitment to its owner, value, $\rho$ and a trapdoor, opened only to
its owner~\cite{shielded}. Their commitments form one append-only Merkle
tree, so a note's path under a root stays valid however the ledger
grows. A leaf's top bit records its kind, set by the chain statement as
it appends, so that a claim can name a note by its position alone.

A spent note is not removed but recorded by its \emph{nullifier} in the
\emph{spent set}, a compact sparse Merkle tree (a key's path is as deep
as its keys' common prefixes: some twenty compressions an access at a
million keys, not 256 levels). The same tree holds, under
domain-separated keys, the state that acknowledgment needs
(Section~\ref{sec:ack}): each pending spend, each locked output, each
acknowledgment and each voided spend, and each recorded anchor's
history. Nothing is ever deleted. A nullifier's leaf names how the note
went: \emph{spent} by a transaction, named by its digest, or
\emph{exited} through contract $c$, by an exit record. A state root
hashes the height, the commitment tree's root and the spent set's root,
and the sequencer signs it (BIP340).

The transaction types are a mint (the hub's), a transfer, a shielded
spend, an acknowledgment, an anchor record (Section~\ref{sec:finality})
and an exit record (Section~\ref{sec:withdraw}). Three proofs are
involved, all RISC Zero (Table~\ref{tab:proofs}).
\begin{table}[!htbp]
\centering
\small
\begin{tabularx}{\linewidth}{>{\raggedright\arraybackslash}p{24mm}>{\raggedright\arraybackslash}p{34mm}L}
\toprule
proof & made by, when & shows; cost (lichen) \\
\midrule
spend & the sender, for each shielded transaction & its notes exist, are hers and are unspent, and the amounts balance; a shield 14.6k cycles, 3.0~s on the laptop; two inputs and two outputs $2^{17}$ to $2^{20}$ cycles, 23 to 110~s \\
chain & the sequencer's prover, for each sealed height & this block is a valid transition, and so was the history before it; 98k cycles (an empty seal) to 0.84M (a payment and its acknowledgment), one segment \\
the hub's statement & the hub, only for a false claim & the claim posted in \tx{request} is false (Section~\ref{sec:statement}); 1.0 to 1.65M cycles, two segments; wrapped in Groth16 for BABE \\
\bottomrule
\end{tabularx}
\caption{The three proofs. The guests are built reproducibly in RISC
Zero's builder image; a contract's statement names the hub's guest's
image id, which names the chain guest's.}
\label{tab:proofs}
\end{table}
No user proves anything for a withdrawal or an exit. A user proves only
her shielded spends, and only on a deployment that has them.

\section{Validity}
\label{sec:validity}

The chain guest proves one sealed height: the block's transactions
applied to the previous state give the new one; its spends are
authorised (a signature, or a verified spend proof for a shielded
note), transparent inputs unlocked, no nullifier spent twice, values
balanced, mints signed by the hub. It takes the previous height's chain
receipt as an assumption, so each receipt vouches for the whole history
from genesis. Its journal carries the height, the state root, the
anchoring state and a history hash
\[
  \mathit{hist}_0 = 0, \qquad
  \mathit{hist}_h = H(\texttt{"lichen/hist"} \,\|\, \mathit{hist}_{h-1} \,\|\, \mathit{SR}_h),
\]
which commits the whole sequence of roots to $h$.

A BIP340 signature, about 260k cycles, is most of a block's work; the
spent set, after it was made compact, about a tenth (lichen). One
machine proving height after height keeps up with a steady load, but
not a burst: on lichen-hub a height took 38 to 157~s, and after a burst
of heights in the trial the prover ran 31 to 35 minutes behind
(lichen). An honest sequencer's lag is only a delay, but the hub's proof
waits for the chain receipt at its claim's state, so the lag counts
against the hub's window (Section~\ref{sec:capacity}). On the RTX 4090 a
height proved in 1.4 to 2.3~s, against 18 to 62~s on the laptop
(lichen).

A valid history is not a unique one: the sequencer can sign two
histories that split, each valid, which the chain proof cannot see. The
wallet and the hub record every root they are shown and refuse to go on
when shown a second root at a height, with the evidence. Anchoring makes
such a split useless.

\section{Finality}
\label{sec:finality}

\paragraph{The seal chain.} The deployment's genesis names an outpoint,
the first \emph{seal}, a P2TR output of the sequencer's key. An
\emph{anchor} spends the current seal and creates the next, beside an
output \tx{OP\_RETURN "lichen"} $\|\; h \;\|\; \mathit{hist}_h$; it is
269~vB (lichen). Bitcoin lets an output be spent once, so there is one
chain of anchors, after Todd's single-use seals~\cite{toddseals}. The
sequencer anchors at most once a Bitcoin block, and only when something
has changed.

\paragraph{Anchors inside the history.} Once an anchor is the margin
deep on Bitcoin (20 blocks on Mutinynet, 6 on mainnet), the sequencer
records it on the L2, and the chain guest checks the record: the
transaction is in the block it names, by Merkle path to the header; it
spends the seal the journal names, whose script is checked so that no
one else could spend it; and it commits a $\mathit{hist}$ that is an
ancestor of the current one. A history forked from the anchored one can
then hold no anchor after the fork.

\paragraph{The contract's work.} The chain guest checks no proof of
work. The hub's statement does: \tx{request}'s block must have headers
on it, each linked and within its own target, whose work sums to at
least $W$, fixed when the contract opens (ten blocks' work at the
difficulty of a checkpoint both sides check against their own view of
the chain). A forger gains nothing from headers declaring low
difficulty: work is summed. This is not a proof of the heaviest chain;
it makes a forged \tx{request} or anchor a matter of mining a fork of
Bitcoin as long as $W$. $W$, the genesis seal, the deployment's keys and
the contract's unit make the deployment's identity,
\[
  \mathcal{D} = H(\texttt{"lichen/contract"} \,\|\, \text{deployment} \,\|\, \text{genesis seal} \,\|\, W \,\|\, u),
\]
which the hub's statement's journal carries (Section~\ref{sec:statement}).

\section{Payments, and the sender's acknowledgment}
\label{sec:ack}

\subsection{The problem}

An exit claim is judged at an anchored state: ``these notes are mine
and unspent at $\Astar$''. The root of each anchor is on Bitcoin, but
the data under it isn't. An operator that anchors a busy block and
serves none of its data leaves every sender in it unable to choose what
to claim:
\begin{itemize}[leftmargin=1.5em]
\item \emph{her inputs:} if the anchor holds her payment, they are
  spent, and a claim of them is false;
\item \emph{her outputs:} if it doesn't, they don't exist, and a claim
  of them is false.
\end{itemize}
Whatever she can't claim defaults to the hub at $\Tc$: the motive. No
better accumulator fixes this: a member's proof needs update data
proportional to the others' changes~\cite{camacho}, which is what an
anchor's data is. The fix is an order: inclusion first, which changes
nothing yet; then the sender's signature, given only once she holds her
data, which makes the spend count.

\subsection{The rules}

\begin{itemize}[leftmargin=1.5em]
\item \emph{Every spend names a one-time key $\mathit{ak}$,} the
  sender's, for that spend only (derived from her wallet's key and the
  spend's first nullifier; for a shielded spend, a public input of its
  proof). Transactions the hub signs, its mints and exit records,
  acknowledge themselves.
\item \emph{A sealed spend is pending.} It spends its inputs and appends
  its outputs, but its outputs are \emph{locked}: no one spends them,
  the sender's change included, until it is acknowledged.
\item \emph{The sender acknowledges against an anchor.} Once an anchor
  at or above the spend's height is recorded, her wallet checks the
  anchor on Bitcoin, the sequencer's state against the roots it signed,
  and that the roots from her block to the anchor hash-chain to the
  $\mathit{hist}_a$ the anchor commits. Then it signs, under
  $\mathit{ak}$, $H(\texttt{"lichen/ack"} \,\|\, \mathit{txid} \,\|\, a
  \,\|\, \mathit{hist}_a)$. Applied by the chain guest (the spend
  pending, not void, the anchor $a$ recorded with that history, $h \le
  a$), it makes the spend count and unlocks its outputs. Naming the
  anchor rather than the spend's own height keeps seals as frequent as
  the sequencer likes.
\item \emph{An exit can void a pending spend.} An exit record retiring
  a nullifier that a pending spend spent voids that spend: its outputs
  locked for good, a later acknowledgment refused. One spent by an
  acknowledged spend makes the record invalid.
\item \emph{The sequencer's policy,} the hub's protection rather than a
  guest rule: it sequences an acknowledgment only once its anchor is
  the margin deep, and not if a \tx{request} on Bitcoin claims the
  spend's inputs. The earlier of an exit and an acknowledgment wins.
\end{itemize}

\subsection{Who is protected}

\begin{center}
\small
\begin{tabularx}{\linewidth}{>{\raggedright\arraybackslash}p{30mm}>{\raggedright\arraybackslash}p{46mm}L}
\toprule
collusion & what they try & what happens \\
\midrule
the operator, against Alice & anchor her payment, withhold its data & she never acknowledges; it counts for nothing, her inputs stay claimable. A stall, not a theft \\
Alice and the operator, against Bob & never sequence her acknowledgment; she claims her inputs & the exit record voids her payment and Bob's note. He never counted it \\
Alice and Bob, against the hub & she gives the acknowledgment to Bob alone, and claims her inputs & the hub pays her; its exit record voids the payment, so Bob's claim of his note is refuted as voided. The hub pays once \\
Alice, against the hub & claims her inputs through one contract and her acknowledged change through another & her acknowledgment, posted in the one, is the hub's evidence against the other \\
\bottomrule
\end{tabularx}
\end{center}

\textbf{A payee who waits for the acknowledgment to be sequenced is safe
against any collusion}: from then no exit can void the spend. Before
that he holds a soft payment, the sequencer's word and then Alice's
acknowledgment.

\subsection{What it costs}

\begin{itemize}[leftmargin=1.5em]
\item \emph{Time.} A payment is firm, and spendable onward, once its
  anchor is the margin deep and the acknowledgment sequenced. On
  Mutinynet, Bob's payment sent just after an anchor's broadcast, the
  worst case of two anchors, was firm in 39 blocks, 20 minutes; Cleo's,
  sent before one, in 21 blocks, 10 minutes (lichen). On mainnet,
  about an hour (est.). There are no fast onward transfers in this
  version.
\item \emph{Presence.} The sender is online until the anchor after her
  payment is the margin deep. If she leaves first, her payment waits,
  counting for nothing, and she acknowledges on her return.
\item \emph{Proving.} About one BIP340 signature in the chain guest; a
  payment and its acknowledgment are 0.84M cycles, one segment (lichen).
\item \emph{On chain.} Nothing in normal operation. An exit claiming a
  locked note carries its acknowledgment: under 20~vB.
\end{itemize}

This is the confirmation signature of Plasma~\cite{plasma} and Plasma
MVP~\cite{plasmamvp}, with three changes: it binds to a history fixed on
Bitcoin, $(a, \mathit{hist}_a)$, rather than to a block; it counts by
existence and claims carry it, so the operator can't suppress it by
censoring; and outputs are locked until it is sequenced, with an exit
voiding a spend still unacknowledged, so that a payee who waits is firm
(Section~\ref{sec:precedents}).

\section{The claim}
\label{sec:claim}

Alice claims a definite set of her notes, and $b$, at most $\vmax$, in
the contract's units. The claim is: \emph{these notes are mine, live at
$\Astar$, and worth at least $b$.} A note is live if its commitment is
in the tree and its nullifier is not spent. The claim carries:
\begin{itemize}[leftmargin=1.5em]
\item each transparent note by its position (4 bytes); each shielded
  note by its position and opening (value, $\rho$ and trapdoor: 76
  bytes), since the hub knows no shielded opening;
\item for each claimed note still locked, the acknowledgment of the
  spend that made it;
\item one signature per owner key on the claim's fingerprint
  \[
    E = H(c \,\|\, O_D \,\|\, b \,\|\, \text{positions} \,\|\, H(\text{acknowledgments}) \,\|\, \ldots),
  \]
  $O_D$ being the contract's deposit coin, which names the contract on
  Bitcoin; a shielded owner's signature is by its exit key, beside the
  keys that derive its nullifiers (Section~\ref{sec:shielded}).
\end{itemize}
\tx{request} posts the claim raw in its witness, at most about 1,200
bytes in items of up to 80 bytes, the most Bitcoin Core relays in a
tapscript witness item, beside Alice's one-time Winternitz signatures on
$b$ and on $E$. Script cannot check the claim's bytes; the hub's
statement checks that they hash to the signed $E$, so a relayer can't
swap them. A claim holds up to 64 notes, 13 of them shielded.

\paragraph{When it is judged: $\Astar$ and the margin.} $\Astar$ is the
last anchor confirmed at least a margin $m$ of blocks before
\tx{request}'s block. The margin is for censorship. If the operator holds
a payment of hers unsequenced, she sees at the tip that it isn't
included and claims its inputs; without a margin, the operator could
include the payment in an anchor confirming before her \tx{request} and
make her claim false. With it, what she saw when she broadcast is what
counts, provided \tx{request} confirms within $m$ blocks, which she
ensures by fee-bumping. The tree only grows and a spend is permanent, so
the one state at $\Astar$ answers every question a claim raises, and the
hub's proof needs one chain receipt.

\paragraph{What the wallet claims.} Its notes in $\Astar$'s tree,
unlocked or carrying their acknowledgment, the largest first, up to the
cap and the witness's room. A claim spends whole notes and the contract
pays once, so value above $b$ is forfeit; to withdraw an exact amount
the wallet first pays itself a note of that size. A note newer than
$\Astar$ isn't in its tree, so once her notes are anchored the wallet
waits until that anchor is the margin deep. With the sequencer gone, it
claims from the notes it kept as it read the L2.

\section{The hub's statement}
\label{sec:statement}

\emph{The claim posted in \tx{request} is false.} The statement is fixed
when the contract opens: the image id of the hub's guest, which names
the chain guest's, and the journal $c \,\|\, \mathcal{D} \,\|\, H(O_D
\,\|\, m)$. It reads \tx{request} from Bitcoin, by its inclusion in a
block with the contract's work on it, and takes $b$ and $E$ from Alice's
signatures' digits: Bitcoin checked those signatures in \tx{request}'s
script, so her keys needn't be in the journal. It holds on any one of
these grounds (Table~\ref{tab:grounds}).
\begin{table}[!htbp]
\centering
\small
\begin{tabularx}{\linewidth}{>{\raggedright\arraybackslash}p{26mm}L}
\toprule
ground & shown by \\
\midrule
malformed & the claim doesn't parse, doesn't hash to her $E$, repeats a note, or $b$ exceeds the cap \\
not hers & a note's owner signed none of the claim's signatures \\
too much & the notes are worth less than $b$ \\
not a note & at $\Astar$, the leaf at a note's position isn't its commitment (an empty leaf past the tree's end included) \\
spent & at $\Astar$, its nullifier is spent by a spend that counts: hub-signed, acknowledged at $\Astar$, or pending with an acknowledgment the hub shows from anywhere (the L2, another claim on Bitcoin) whose anchor confirmed before the cutoff. This contract's own exit record does not count \\
locked & at $\Astar$, it is an output of a pending spend, and the claim carries no valid acknowledgment of it \\
voided & at $\Astar$, it is an output of a spend an exit record voided \\
a second claim & the note's owner signed a claim of it through another contract \\
\bottomrule
\end{tabularx}
\caption{The grounds on which the hub's statement holds.}
\label{tab:grounds}
\end{table}

\paragraph{$\Astar$ is the last.} The hub shows $\Astar$'s anchor
confirmed by the cutoff, and the transaction spending $\Astar$'s seal
confirmed after it. Only the sequencer's key spends the seal, and Bitcoin
allows one spend, so this is objective. If nothing has spent the seal
yet, the hub makes the sequencer anchor (on Mutinynet it minted itself
a sat, so that there was something to seal), which it can always do. The
state at $\Astar$ comes from the sequencer's chain receipt at the
height where $\Astar$'s own anchor record is applied: anchor records are
applied in seal order, so every anchor up to $\Astar$ is there, each
confirmed before the cutoff.

\paragraph{Why the operator can't forge a refutation.} The state comes
from a chain receipt, so it is valid; it is the anchored history, and
the anchor and \tx{request} carry the contract's work. A nullifier
enters the spent set only by a spend, which needs the note's keys, or by
an exit record, which the chain guest accepts only with the owners'
signatures on that claim. A spend counts against Alice only with her
acknowledgment, which she gives only with her data. A second claim needs
her signature, and a copied signature names its own contract. So every
refutation traces to something the owner did.

\paragraph{The two kinds of honesty.} For an honest claimant no
refutation exists: her notes are in the tree at $\Astar$, the spends
she hasn't acknowledged don't count, and the spends she has
acknowledged she has left out of her claim. For an honest hub, every
false claim has one: whatever makes the claim false is in the
operator's data. The statement's guest is the decision the hub makes
natively, in milliseconds, so the hub proves only what it has already
found to hold.

\paragraph{Two claims of one note.} Requested together through two
contracts, each is the other's second claim, and both are asserted
false: a double claimer loses both contracts. After deciding a claim
holds, the hub keeps looking for a second claim of its notes until a
quarter of its window is left.

\section{Shielded notes}
\label{sec:shielded}

An exit must claim shielded notes directly: unshielding is an L2
transaction, which needs the sequencer, which an exit can't count on.
The hub knows no shielded opening, and to check a note unspent it needs
the nullifier, which needs the owner's key. So the keys are split, and a
trapdoor $r$, chosen by the sender and encrypted to the receiver with
the value, enters the commitment and the nullifier:
\begin{align*}
  \mathit{owner} &= H(\texttt{"lichen/owner"} \,\|\, H(\mathit{ask}) \,\|\, \mathit{nk} \,\|\, X),\\
  \mathit{cm} &= H(\texttt{"lichen/snote"} \,\|\, \mathit{owner} \,\|\, v \,\|\, \rho \,\|\, r),\\
  \mathit{nf} &= H(\texttt{"lichen/snf"} \,\|\, \mathit{nk} \,\|\, \rho \,\|\, r).
\end{align*}
$\mathit{ask}$ authorises spends inside the spend proof and is never
revealed. A claim reveals $H(\mathit{ask})$, $\mathit{nk}$ and $X$ once
per owner, and each note's opening; $X$, a Schnorr key, signs $E$ and
nothing else. The hub then checks everything natively, on the same
grounds as for transparent notes, and ``not hers'' becomes: the owner
isn't $H(H(\mathit{ask}) \,\|\, \mathit{nk} \,\|\, X)$ of what was
posted, or $X$'s signature doesn't verify.

Revealing $\mathit{nk}$ lets anyone link the claimed notes' nullifiers
to them, and lets a sender who knows a note's $\rho$ and $r$ see when
the owner spent it; nobody can spend with what a claim reveals. So
after an exit the wallet moves to a fresh shielded address. The
trapdoor also made the commitment hiding: before it, anyone with an
address could test leaves against $H(\mathit{owner} \,\|\, v \,\|\,
\rho)$ for every plausible value $v$, $\rho$ being public. The
alternatives were key-image nullifiers, which put elliptic-curve
multiplications in every shielded spend proof (from $2^{17}$ to about
$2^{20}$ cycles, est.), a cost on every payment for an exit that should
be rare; and a proof by Alice, which is what this design removes.

On Mutinynet, Sara shielded 90,000 sats (3.0~s on the laptop), then
exited claiming her shielded note and a transparent one: her
\tx{request} was 2,194~vB, 59 more than a one-note transparent claim,
and \tx{request\_timeout} paid her (lichen).

\section{Withdrawals}
\label{sec:withdraw}

\paragraph{Cooperative.} Alice's wallet builds the same claim an exit
would post, and the fold from the contract's terms, and sends both to
the hub. The hub checks the claim against its own ledger, now rather
than at an anchor (each note hers, unspent, unlocked or carrying its
acknowledgment, the owners' signatures for this contract, the value at
least $b$), sequences its \emph{exit record}, which retires the notes
and, being hub-signed, counts at once, and then signs and broadcasts the
fold. There is no wait for an anchor or a proof: if the hub records the
claim and won't fold, her exit posts the very same claim, which this
contract's own record doesn't refute; if it never recorded it, the notes
are live. Either way she is paid, so she needn't know which. On
Mutinynet the hub's part took 123~ms and the fold was 205~vB; Olga's
withdrawal took 20 minutes end to end, all of it waiting to acknowledge
the split that sized her note (lichen).

\paragraph{The exit.} If the hub won't fold, or isn't there, she posts
the claim in \tx{request}. As soon as the hub sees it, it holds the
claimed notes: the sequencer refuses their spends and any acknowledgment
of a spend of them. It decides the statement natively. If the claim
holds, it asserts nothing, records the exit on the L2, and lets
\tx{request\_timeout} pay her after $w_H$. If it is false, it proves and
asserts, and her wallet checks the proof natively (about 1.5~ms): a
valid one she concedes, an invalid one she challenges. The transactions
and BABE's part are in the main paper~\cite[App.~A]{lichen}.

\section{Measurements}
\label{sec:measured}

All on Mutinynet (lichen), seven contracts of $D = 100{,}000$ sats with
100,000 of liquidity: $\vmax = 199{,}155$, 255 units of 781 sats, $r_A
= 20{,}000$, $r_H = 31{,}816$, windows $w_H = 240$ blocks (2 hours),
$\Delta_2 = 2{,}880$ (a day), $\Delta_1 = 240$.

\subsection{Opening}

\begin{center}
\small
\begin{tabularx}{\linewidth}{>{\raggedright\arraybackslash}p{50mm}L}
\toprule
on chain & 517~vB: the deposit (212) and the funding (305) \\
Alice's device & 78 instances garbled in 9.3~s on the laptop; 10 kept uploaded, 23.9~MiB each, 239~MiB, 65~s on a home connection (some 30~Mbit/s) \\
the hub & 20 to 36~s checking the 68 opened; 239~MiB stored per contract \\
signatures & Alice's on 17 pre-signed transactions, the hub's on 7, and its 1,016 adaptor pre-signatures \\
cut-and-choose & $\binom{78}{10}^{-1} = 2^{-40.2}$ \\
\bottomrule
\end{tabularx}
\end{center}
The upload is what the swapped roles cost: under the earlier design the
hub garbled and Alice downloaded. Six wallets opening at once over one
home uplink took about half an hour.

\subsection{The exits}

\begin{table}[!htbp]
\centering
\small
\begin{tabularx}{\linewidth}{>{\raggedright\arraybackslash}p{14mm}>{\raggedright\arraybackslash}p{36mm}L>{\raggedright\arraybackslash}p{17mm}}
\toprule
contract & the claim & how it ended & on chain \\
\midrule
1002 & true & the hub found no ground; \tx{request\_timeout} paid her after 2 hours & 2.7~kvB \\
1005 & true, a shielded note and a transparent one & the same & 2.8~kvB \\
1003 & stale: its note paid away and acknowledged & proved false on lichen-hub, asserted 54 minutes after \tx{request}; conceded & 9.9~kvB, then \tx{success} \\
1001 & stale, as 1003 & asserted; challenged anyway (a trial switch); answered in two blocks & 30.8~kvB \\
1004 & true, against a hub restarted hostile & a forged proof asserted; challenged; no instance opened; \tx{failure} paid her & 31.2~kvB \\
\bottomrule
\end{tabularx}
\caption{Five exits on Mutinynet. The other two contracts were
withdrawn cooperatively.}
\label{tab:exits}
\end{table}

\paragraph{Sizes.} \tx{request} 2,135 to 2,194~vB, its floor paid out of
the contract; \tx{request\_timeout} 565~vB; \tx{assert} 7,725~vB;
\tx{challenge\_assert} 20,717~vB; \tx{wrongly\_challenged} 262~vB;
\tx{failure} 589~vB. Sweeping the bit outputs, each of which reads
Alice's signature on $b$ again, is about 2.3~kvB, paid from her fee
budget. The three links' floors, 30,826 sats, and their three 330-sat
anchors make 31,816: exactly $r_H$, which the hub sizes for the deepest
dispute. In contract 1004, $r_A + r_H$ less that left Alice exactly her
own reserve, besides her $b$.

\paragraph{The hub's proving.} The hub's statement for a stale claim
proved in 520.6 and 526.0~s on lichen-hub (1.62M and 1.65M cycles, two
segments) and wrapped in Groth16 in 128.4 and 119.3~s. Of the 54
minutes from \tx{request} to \tx{assert} in contract 1003, 35 waited
for the chain prover, behind after a burst of heights, and 11 were
proving and wrapping. On regtest the same kind of statement (1.02M
cycles) proved in 827~s on the 4-vCPU server and wrapped in 268~s; on
the RTX 4090 it proved in 3.5~s (lichen). RISC Zero's GPU Groth16
prover failed on that card each time (an illegal memory access), so the
wrap stayed on the CPU, about 4.5 minutes: nearly all of a GPU hub's
proving time. With a working GPU wrap the whole proof would be some 10
to 20~s (est.).

\paragraph{Times.} A true claim was paid 241 blocks after
\tx{request}, about 2 hours, and swept 41 blocks later, with no prover on
Alice's side and nothing needed from the hub or the sequencer after her
\tx{request}. Against the lying hub, contract 1004 was paid 274 blocks, 2 hours 23
minutes, after its \tx{request}.

\subsection{Regtest}

The cases not played on Mutinynet ran on regtest, on the mock
statement: one note claimed through two contracts at once (both
asserted false, the double claimer losing both); an exit after a
withdrawal the hub recorded but didn't fold (paid); the hub and
sequencer gone (paid by timeout from the notes kept); a payment never
acknowledged with its input claimed (paid; the payment voided and the
payee's note locked for good); and a payment acknowledged whose
acknowledgment the operator withheld (the locked change claimed,
carrying the acknowledgment: paid, the payee's note unlocked).

\section{What the earlier design had, and this one doesn't}
\label{sec:gone}

\begin{itemize}[leftmargin=1.5em]
\item \emph{Alice's proof,} and her prover: no STARK, no Groth16 wrap,
  no x86 machine with Docker and 8~GB, for any withdrawal or exit.
\item \emph{The counter,} and Alice's worst case: proving a forged
  counter false meant verifying a Groth16 proof inside RISC Zero, 56.3M
  cycles, about an hour on a laptop. The hub now carries the only
  proof, from data it holds.
\item \emph{The return,} a payment of $b$ to the hub with the contract's
  memo, its withdrawal statement and its own dispute; and the
  \tx{respond} link, in which the hub signed $b$. A withdrawal is now a
  claim, the same object cooperatively or on chain.
\item \emph{Exit kits} and the exit set's chain receipt: Alice needs her
  notes' positions and openings, which her wallet has.
\item \emph{The keyless tower.} Under the swapped roles a challenge
  reveals Alice's secrets, so she challenges herself, in a window of a
  day on Mutinynet (Section~\ref{sec:open}).
\end{itemize}

\section{Precedents}
\label{sec:precedents}

How other designs treat a hostile operator, and where lichen differs.
\begin{center}
\small
\begin{tabularx}{\linewidth}{>{\raggedright\arraybackslash}p{26mm}>{\raggedright\arraybackslash}p{38mm}>{\raggedright\arraybackslash}p{30mm}L}
\toprule
& when data is withheld & sender sign-off & what L1 enforces \\
\midrule
Plasma~\cite{plasma} & mass exit, coordinated & both parties acknowledge after inclusion & bonds, challenge games \\
Plasma MVP~\cite{plasmamvp} & exit within 7 days & confirmation signatures & a priority queue \\
More Viable Plasma~\cite{morevp} & in-flight exits, canonicity games & dropped & youngest-input priority \\
Plasma revisited, 2023~\cite{neoplasma} & forced inclusion, or L1 reverts & none & a queue, reversion \\
Intmax2~\cite{intmax2} & unsigned counts for nothing & BLS-signs the root before posting & the signer set, balance proofs \\
lichen & she claims; the hub proves the claim false or pays & signs $(\mathit{txid}, a, \mathit{hist}_a)$; outputs locked until then & nothing about the L2: an anchor is an \tx{OP\_RETURN} \\
\bottomrule
\end{tabularx}
\end{center}
\begin{enumerate}[leftmargin=1.5em]
\item \emph{Bitcoin enforces nothing about the L2,} so every rule is
  applied after the fact, in a dispute. The precedents use L1 to stop a
  bad state or rescue funds; Bitcoin can't reject an anchor, freeze, or
  revert. Lichen never prevents a bad anchor: the hub's statement
  decides which anchored facts count against a claim.
\item \emph{Per-user collateral, not a pool.} A Plasma exit draws on one
  contract, so a double exit steals from everyone, hence priorities and
  coordination. Here each exit is against one contract's collateral: no
  queue, no ordering among users.
\item \emph{The burden on the data holder.} Everywhere else the exiter
  proves something: inclusion, a path, a balance proof. Here she claims
  and the operator proves, since it holds all the data and pays the
  exits.
\item \emph{No in-flight exits.} MoreVP's in-flight exit claims ``the
  inputs or the outputs, whichever is canonical'', with bonded games of
  several parties. On Bitcoin each party in a game is another
  garbled-circuit setup; acknowledgment keeps claims definite and
  disputes two-party. MoreVP's objections to confirmation signatures
  still apply and are accepted: the extra step, griefing by withholding
  (a stall), and the sender's free option until she signs, which the
  payee closes by waiting.
\end{enumerate}

\section{Limits and open questions}
\label{sec:open}

\subsection{The hub's proving capacity}
\label{sec:capacity}

The hub must prove every false claim within $w_H$, and each one it fails
to answer is paid. A liar loses a contract, so false claims are costly
one at a time; but Sybils who open many contracts, pay their tokens to
each other and exit together could outrun one prover. One proof per
contract is needed, since each journal is its contract's. On lichen-hub
a proof took about 11 minutes with its wrap, after the chain prover had
caught up; on a GPU the proving is seconds and the wrap the rest.
Options: a batched STARK with a small composing guest and a wrap per
contract; a working GPU wrap; a longer $w_H$, which is also the honest
user's wait; a limit on contracts opened per window.

\subsection{Block space}

Each exit is a dispute competing for block space. Against a hub that
asserts falsely against everyone, each user's challenge is 20.7~kvB
within $\Delta_2$; with $N$ contracts and a window of $w$ blocks,
$N \cdot 20.7\,\mathrm{kvB} \le w \cdot 1\,\mathrm{MvB}$ allows about
7,000 contracts a day-long window on mainnet (est.), more as the window
grows. Each false assertion also costs the hub 7.7~kvB and the
contract. Against a wave of false claims, the honest hub's 7.7~kvB
assertions must fit within $w_H$.

\subsection{Data availability}

Acknowledgment closes the data-availability gap for funds: withholding
stalls payments but takes nothing. What remains are assumptions: the
operator can stop anchoring, or not sequence acknowledgments, so that
payments stall; and softer finality (spends of locked outputs that count
only if their parents do, as with unconfirmed Bitcoin outputs, or a
hub-signed receipt for an acknowledgment) needs analysis before it is
built.

\subsection{The rest}

\begin{itemize}[leftmargin=1.5em]
\item \emph{A shielded input's lock} is kept by the wallet and the
  sequencer, not the chain guest, since its position is private; the
  guest should enforce it.
\item \emph{A tower for the challenge.} A threshold of towers holding
  Alice's secrets, so that only a quorum colluding with the hub could
  steal; or accept that challenging is hers, in a window of days.
\item \emph{The upload,} 239~MiB per contract: fewer kept instances at
  larger $N$, or Duty-Free Bits' smaller encoding~\cite[App.~A]{lichen}.
\item \emph{The seal's other spenders.} If the seal gains a freeze path,
  the first spend of $\Astar$'s seal might be the freeze; the
  last-anchor argument needs only some spend after the cutoff, but this
  should be checked.
\item \emph{The security argument,} here a sketch: honest-claimant
  soundness (no refutation of a true claim) and the hub's completeness
  (a refutation of every false one) written as proofs, beside BABE's
  own and the cut-and-choose's.
\end{itemize}

\begin{thebibliography}{99}

\bibitem{lichen} waxwing. Lichen: L2s in channels enforced by Bitcoin. Version 0.5, November 2026. \url{https://github.com/AdamISZ/lichen/blob/master/papers/lichen-v0.5.0.pdf}

\bibitem{lichenbuilt} waxwing. Lichen as built: the BABE gadget, the hub's ledger and the exit. Version 0.1 (not maintained), October 2026. \url{https://github.com/AdamISZ/lichen/blob/master/papers/lichen-built-v0.1.1.pdf}

\bibitem{lichencode} lichen: the ledger, the sequencer, the hub and the wallet. \url{https://github.com/AdamISZ/lichen}

\bibitem{explainers} lichen's explainers: opening a contract, payments and acknowledgments, a cooperative withdrawal, exiting against a hostile operator, the payout, the proving pipeline. November 2026. \url{https://github.com/AdamISZ/lichen/tree/master/docs/explainers}

\bibitem{shielded} lichen's shielded transfers: an example L2. \url{https://github.com/AdamISZ/lichen/blob/master/docs/SHIELDED.md}

\bibitem{toddseals} P.~Todd. Preventing consensus fraud with commitments and single-use-seals, 2016. \url{https://petertodd.org/2016/commitments-and-single-use-seals}

\bibitem{camacho} P.~Camacho and A.~Hevia. On the impossibility of batch update for cryptographic accumulators. LATINCRYPT 2010.

\bibitem{plasma} J.~Poon and V.~Buterin. Plasma: scalable autonomous smart contracts, 2017. \url{https://www.plasma.io/plasma.pdf}

\bibitem{plasmamvp} V.~Buterin. Minimal viable Plasma, 2018. \url{https://ethresear.ch/t/minimal-viable-plasma/426}

\bibitem{morevp} More Viable Plasma, 2018. \url{https://ethresear.ch/t/more-viable-plasma/2160}

\bibitem{neoplasma} V.~Buterin. Exit games for EVM validiums: the return of Plasma, 2023. \url{https://vitalik.eth.limo/general/2023/11/14/neoplasma.html}

\bibitem{intmax2} E.~Rybakken, L.~Hioki and M.~Yaksetig. Intmax2: a ZK-rollup with minimal onchain data and computation costs featuring decentralized aggregators. ePrint 2023/1082. \url{https://eprint.iacr.org/2023/1082}

\end{thebibliography}

\end{document}
